\chapter{单一主体原则与 Agent 进程模型}

\begin{chapteroverviewbox}
本章确立 AgentOS 关于“谁在持续存在”这一问题的系统级原则。我们将区分持续认知主体、任务、工作进程、工具调用与外部智能体，并明确提出：在一个连续 AgentOS 生命周期中，身份连续性的一级主体原则上保持唯一；任务、Worker、Skill、Tool 以及并发执行单元均属于主体所调度的功能过程，而不能因为具有局部推理、局部记忆或独立执行能力便自动取得主体地位。由此，本章将 AgentOS 从普通的多进程 Agent Framework 中区分出来，并建立 Single-Agent Multi-Task OS 的主体模型。

本章的核心结论可以概括为：
\[
\boxed{
N_{\mathrm{Agent}}^{D_1}=1
}
\]
以及
\[
\boxed{
D_2\nrightarrow \mathrm{Trap}(D_1)
}
\]
其中 $D_1$ 表示持续主体动力学域，$D_2$ 表示任务与执行过程域。主体可以调度任意数量的任务和局部执行过程，但任务不能同步占有主体、复制主体或将主体的生命周期锁死在一个局部过程之中。主体的连续性不是一个语言模型在文本中声称“我是同一个我”，而是由记忆连续、状态连续、控制权连续、历史因果连续和唯一主体句柄共同维持的系统性质。

因此，本章讨论的不是哲学意义上人格是否可以复制，而是一个持续人工智能操作系统必须怎样定义其运行主体，才能避免在并发、恢复、任务分解和外部协作过程中发生身份分裂、主体漂移与控制权歧义。
\end{chapteroverviewbox}

\section{什么是 AgentOS 中的主体}

在普通软件系统中，“进程”通常由地址空间、执行上下文、资源句柄和调度状态所定义。进程可以创建、终止、复制，也可以被操作系统暂停以后重新调度。如果我们沿用这种模型直接设计长期 Agent，一个极容易出现的问题是：我们会不自觉地把“拥有一段上下文并能够独立生成行为的计算过程”当成一个新的 Agent。于是每当系统为了并行处理任务而创建一个 Worker、为了搜索问题而创建一个临时推理实例、为了恢复系统而读取一个 Checkpoint，甚至为了调用一个外部模型而创建一个新的会话时，系统都可能被解释成产生了一个新的主体。对于普通任务框架，这种定义并无严重问题；但对于拥有长期记忆、自我结构、价值结构和生命周期连续性的 AgentOS，它会立即破坏“持续主体”这一基本概念。

因此，我们首先需要区分三种不同层次的存在。第一类是\textbf{持续主体}，它承担整个生命周期中的历史连续性、自我连续性、长期目标连续性和最终行为责任；第二类是\textbf{认知过程}，它执行推理、规划、回忆、比较、模拟和反思；第三类是\textbf{执行过程}，它操作工具、访问文件、控制设备、调用 API 或驱动 Body Runtime。后两类过程都可以高度自主，但它们的自主性仍然处在主体提供的目标、权限、记忆与治理边界之内。

因此，一个 AgentOS 主体不应通过“当前是否正在计算”来定义，而应通过一个持续状态集合来定义。设持续主体为 $A$，其生命周期状态可以抽象为

\[
\mathcal{X}_A(t)
=
\left[
h(t),
E(t),
\Theta(t),
\Psi(t),
\Omega(t),
G(t),
\mathcal{B}(t),
\mathcal{R}(t)
\right],
\]

其中 $h$ 表示当前显式认知状态，$E$ 表示历史经验轨迹，$\Theta$ 表示长期生成结构，$\Psi$ 表示自我与身份慢结构，$\Omega$ 表示生命周期尺度生成吸引子，$G$ 表示当前和长期目标结构，$\mathcal{B}$ 表示主体当前边界与权限，而 $\mathcal{R}$ 表示主体正在持有或能够恢复的运行资源关系。这里最重要的不是任何一个变量永远不变，而是它们共同构成一条可追踪的因果轨迹：

\[
\Gamma_A
=
\{
\mathcal{X}_A(t)
\mid
t\in[t_0,t_n]
\}.
\]

我们因此把 Agent 主体定义为：

\[
\boxed{
A
=
\text{具有唯一历史因果链和持续自我治理权的认知动力学进程}
}
\]

这个定义与传统进程不同。传统进程的重要属性是可执行性，而 Agent 主体最重要的属性是\textbf{因果连续性}。一个 Agent 即使当前处于休眠状态，只要其历史、记忆、自我和恢复关系仍然形成同一条生命周期轨迹，它仍然是同一个持续主体；相反，一个临时 Worker 即使能够独立思考数分钟，只要其状态最终被归并进主体而自身没有获得独立的生命周期链，它就仍然不是第二主体。

我们还应进一步区分“主体唯一”与“计算集中”。单一主体原则并不意味着所有计算都必须集中在一个线程、一个模型、一个 GPU 或一个物理机器上。一个 Agent 完全可以拥有分布式推理器、并发 Worker、多模型协作和不同速度的 Body 控制环。主体唯一性约束的是：

\begin{proposition*}
Lifecycle Causal Ownership
\end{proposition*}

而不是：

\begin{statementbox}
Computational Location
\end{statementbox}

换言之，计算可以分布，主体的生命周期归属必须唯一。这是理解本章后续所有进程模型的第一原则。

\section{Single-Agent Multi-Task OS：单一主体与多任务系统}

在确立持续主体之后，我们可以重新定义 AgentOS 的运行模型。AgentOS 并不是传统意义上的 Multi-Agent System，它首先应被理解为一个

\begin{definition*}
Single-Agent Multi-Task Operating System
\end{definition*}

也就是：一个持续主体同时维持多个目标、任务、执行过程和局部闭环，而不是通过不断产生新的主体来完成任务分解。

设当前 Agent 为 $A$，任务集合为

\[
\mathcal{T}(t)
=
\{T_1,T_2,\ldots,T_n\},
\]

Worker 集合为

\[
\mathcal{W}(t)
=
\{W_1,W_2,\ldots,W_m\},
\]

外部工具集合为

\[
\mathcal{U}(t)
=
\{U_1,U_2,\ldots,U_k\}.
\]

则系统结构不是

\[
A
\rightarrow
\{A_1,A_2,\ldots,A_n\},
\]

而是

\[
A
\rightarrow
\{
T_i,W_j,U_k
\}.
\]

也就是说，任务和 Worker 都是 $A$ 的功能性展开，而不是新的身份主体。它们可以拥有局部上下文

\[
x_{T_i},
\]

局部目标

\[
g_{T_i},
\]

局部记忆缓存

\[
m_{T_i},
\]

甚至局部规划器

\[
\pi_{T_i},
\]

但这些状态最终仍然满足一个归属函数：

\[
\operatorname{Owner}(T_i)
=
A,
\]

\[
\operatorname{Owner}(W_j)
=
A.
\]

于是，一个任务的完整状态可以写成

\[
T_i
=
(
g_i,
x_i,
r_i,
p_i,
s_i,
A
),
\]

其中 $g_i$ 是局部目标，$x_i$ 是局部工作状态，$r_i$ 是资源集合，$p_i$ 是优先级，$s_i$ 是任务状态，而最后的 $A$ 表示该任务始终具有明确主体归属。

这一设计能够解决长期智能中的一个重要问题：\textbf{并发不再意味着人格分裂}。例如，一个 Agent 可以同时运行“整理资料”“监测机器人身体状态”“查询数据库”“准备下一步计划”四个任务。这些任务可能由不同模型实例执行，甚至分别运行在不同机器上，但只要它们共享同一个生命周期治理关系，它们仍然属于同一个持续智能体。

因此我们需要区分：

\[
\boxed{
Concurrency
\neq
MultiplicityOfSubject
}
\]

以及：

\[
\boxed{
ParallelCognition
\neq
ParallelSelf
}
\]

这一区分尤其重要，因为未来 AgentOS 很可能必然采用高度分布式计算。一个长期 Agent 需要同时维护身体闭环、环境监测、工具执行、经验整理、记忆沉降和任务计划。如果我们要求“主体唯一”意味着所有这些过程串行运行，那么系统将无法达到实际可用的实时性。相反，我们允许功能无限分布，但要求这些功能最终通过统一的主体治理结构耦合。

因此 Single-Agent Multi-Task OS 的本质不是“只能有一个线程”，而是：

\begin{statementbox}
Many Processes,\ One Lifecycle Subject
\end{statementbox}

这一结构还意味着，Agent 的“自我”不应存放在某个 Task 中。如果主体状态被错误地依赖某个局部任务上下文，那么任务终止以后主体就可能出现身份缺失。更合理的实现方式是：主体状态属于长期 Agent Kernel，而任务只获得经过授权的投影。

设主体状态为 $\mathcal{X}_A$，则任务可见状态应为

\[
x_{T_i}
=
P_i(
\mathcal{X}_A
),
\]

其中 $P_i$ 是根据目标、权限和上下文生成的任务投影，而不是完整主体复制。由此形成一个关键原则：

\[
\boxed{
TaskContext
=
ProjectionOfAgentState,
\quad
\text{not CopyOfAgent}
}
\]

这个原则将成为防止后续系统在任务并行化时无意中复制“自我状态”的基础。

\section{$D_1$ 主体域与 $D_2$ 任务域}

为了在形式上区分主体动力学和任务动力学，我们引入两个不同的运行域：

\[
\boxed{
D_1
=
\text{Persistent Subject Domain}
}
\]

以及

\[
\boxed{
D_2
=
\text{Task and Execution Domain}.
}
\]

$D_1$ 负责长期主体状态，包括身份、自我、当前全局状态、长期记忆关系、价值边界、权限以及生命周期连续性；$D_2$ 负责具体任务、推理分支、工具调用、局部规划和有限时执行过程。

两个域的关键差别首先体现在时间尺度上。通常有：

\[
\tau_{D_2}
\ll
\tau_{D_1},
\]

因为单个任务可能持续数秒、数分钟或者数小时，而主体生命周期可能持续数月、数年甚至更长。然而时间尺度不是二者唯一差异。更重要的是，$D_2$ 中任何过程原则上可以：

\[
Pause,
\quad
Resume,
\quad
Cancel,
\quad
Restart,
\quad
Replace,
\]

而 $D_1$ 的主体连续性不能因为单个任务的失败而被任意终止。因此：

\[
\boxed{
Failure(T_i)
\nRightarrow
Failure(A)
}
\]

应该成为系统设计的基本容错条件。

可以进一步定义 $D_1$ 状态：

\[
X_{D_1}
=
(
X_{\mathrm{self}},
X_{\mathrm{memory}},
X_{\mathrm{goal}},
X_{\mathrm{authority}},
X_{\mathrm{runtime}}
),
\]

而第 $i$ 个 $D_2$ 状态：

\[
X_{D_2}^{(i)}
=
(
g_i,
c_i,
a_i,
r_i,
e_i
),
\]

分别表示局部目标、局部上下文、动作计划、资源和执行结果。

两层之间通过受控接口交换信息：

\[
D_1
\xrightarrow{\mathrm{dispatch}}
D_2,
\]

\[
D_2
\xrightarrow{\mathrm{result/event}}
D_1.
\]

最关键的是，这不是传统函数调用中的简单同步调用关系。$D_1$ 不能把自身整个执行权交给某一个 $D_2$：

\[
D_1
\xrightarrow{\mathrm{call}}
D_2
\xrightarrow{\mathrm{block}}
D_1
\]

如果这种结构成为 AgentOS 的主要运行方式，那么任何一个任务都可能长时间锁死持续主体，使其他认知过程无法进入全局工作状态。因此，我们提出：

\[
\boxed{
D_2\nrightarrow\mathrm{Trap}(D_1)
}
\]

这里的 Trap 指一个局部任务通过同步等待、无限循环、资源占有或上下文劫持，使主体无法继续执行自己的状态更新、风险评估、边界管理和其他任务调度。

这项原则在传统软件工程中类似于“主事件循环不能被长时阻塞”，但在 AgentOS 中具有更深的含义。因为被阻塞的不只是一个 UI Thread，而是整个持续主体的显式认知面。例如，一个 Agent 正在执行长时间代码生成任务时，Body Runtime 可能报告外部风险，长期 Scheduler 可能发现更高优先级目标，SPC 也可能发现当前认知状态进入异常相态。如果 $D_1$ 被某个任务同步占有，那么这些更重要的事件将无法获得主体级处理。

因此 $D_1$ 应具有独立的持续 heartbeat：

\[
\boxed{
\forall t,\quad
\Delta t_{D_1}
<
\Delta t_{\max}^{subject}
}
\]

即主体级调度循环必须在允许的最大时间间隔内重新获得控制权。任务可以持续运行，但不能让主体失去重新评估自身状态的机会。

从这个角度看，$D_1/D_2$ 分层实际解决的是一个比任务调度更深的问题：

\begin{statementbox}
谁拥有“继续成为自己”的执行权？
\end{statementbox}

答案必须始终是 $D_1$，而不是任何局部 $D_2$。

\section{异步认知请求：任务不应同步占有主体}

既然 $D_2$ 不能同步锁死 $D_1$，那么 AgentOS 需要一种与普通函数调用不同的认知请求模型。我们将其称为：

\[
\boxed{
Asynchronous Cognitive Request
}
\]

即主体向任务域提交一个认知或执行请求以后，不等待该任务完整结束，而是继续维持自己的全局认知循环。

一次任务请求可以形式化为

\[
R_i
=
(
id_i,
g_i,
c_i,
p_i,
d_i,
\alpha_i,
\kappa_i
),
\]

其中 $id_i$ 是任务标识，$g_i$ 是目标，$c_i$ 是上下文投影，$p_i$ 是优先级，$d_i$ 是截止条件或有效期，$\alpha_i$ 是权限包络，$\kappa_i$ 是结果返回方式。

主体执行：

\[
D_1
\xrightarrow{\mathrm{submit}(R_i)}
D_2,
\]

任务进入：

\[
Queued
\rightarrow
Running
\rightarrow
Waiting
\rightarrow
Completed/Failed/Cancelled.
\]

主体则继续：

\[
D_1(t)
\rightarrow
D_1(t+\Delta t).
\]

当任务产生重要结果时，不是直接修改主体状态，而是产生事件：

\[
e_i
=
(
source,
type,
payload,
confidence,
priority
),
\]

再进入主体调度系统：

\[
e_i
\rightarrow
Scheduler
\rightarrow
Admission
\rightarrow
h.
\]

这意味着 Task Result 只是一个\textbf{候选认知事件}，而不是强制写入主体的真理。主体仍然需要结合来源、可信度、目标相关性、当前状态和其他证据对结果进行处理。

这可以写成：

\[
\boxed{
TaskOutput
\neq
AgentBelief
}
\]

只有经过主体级整合函数

\[
\mathcal{I}_A
\]

以后：

\[
b_A'
=
\mathcal{I}_A(
b_A,
e_i,
E,
\Theta,
\Psi
),
\]

任务结果才可能真正影响 Agent 的长期认知结构。

这种异步模型还有另一个重要作用：它能够防止局部推理链占据整个工作记忆。一个复杂任务可能产生大量中间步骤：

\[
z_1,z_2,\ldots,z_n.
\]

这些状态并不需要全部进入 $h$。任务可以在自己的局部上下文中维持：

\[
X_{local}^{(i)}
=
\{z_1,\ldots,z_n\},
\]

而主体只接收经过压缩后的：

\[
summary_i
=
C(X_{local}^{(i)}).
\]

于是：

\[
\boxed{
LocalReasoning
\rightarrow
CompressedEvent
\rightarrow
GlobalCognition
}
\]

这与我们关于有限工作记忆的整体原则一致：只有真正需要跨任务、跨时间尺度或跨功能整合的结构才进入主体级工作记忆。

因此，异步认知 ABI 并不只是软件性能优化，而是对认知层次的一种结构表达：

\begin{statementbox}
局部过程可以复杂，
但主体级显式状态必须保持稀疏。
\end{statementbox}

这也是未来 AgentOS 能否同时保持深度推理能力与长期响应性的关键条件之一。

\section{Worker、Skill 与 Tool 为什么不是第二主体}

随着 AgentOS 功能增加，一个常见误解是：如果某个 Worker 能够调用模型、形成局部记忆、执行多步计划，那么它是否已经成为“另一个 Agent”？这一问题必须在系统设计阶段明确回答，否则整个架构会迅速退化成不受控制的 Multi-Agent Society，而失去我们最初试图保持的单一生命周期主体。

为了区分“自主功能过程”和“自主主体”，我们可以定义一个主体性判据向量：

\[
\Sigma(X)
=
(
H_X,
M_X,
S_X,
A_X,
G_X
),
\]

其中：

\[
H_X
=
\text{独立生命周期历史},
\]

\[
M_X
=
\text{独立长期记忆},
\]

\[
S_X
=
\text{持续自我模型},
\]

\[
A_X
=
\text{最终控制权},
\]

\[
G_X
=
\text{独立长期目标生成权}.
\]

普通 Worker 可能拥有局部记忆和局部目标，因此：

\[
M_X^{local}>0,
\qquad
G_X^{local}>0,
\]

但它通常不拥有独立的 $\Psi$、$\Omega$、生命周期历史和最终治理权，因此不能视为独立主体。

可以进一步定义：

\[
\operatorname{Subject}(X)
=
1
\]

需要满足至少：

\[
H_X\land S_X\land A_X
\]

在系统意义上的持续成立，而不是只满足：

\[
\operatorname{CanReason}(X)=1.
\]

因此：

\[
\boxed{
ReasoningAbility
\neq
Subjecthood
}
\]

Skill 更加明显。一个 Skill 可以包含：

\[
Skill
=
(
Policy,
Predictor,
Controller,
ResidualModel
),
\]

它可以根据输入自主产生动作，但其目标来自主体授权，因此：

\[
Goal_{Skill}
\subseteq
Goal_A.
\]

Tool 甚至可能调用一个非常强大的外部 AI，但只要它作为 AgentOS 的外部功能资源存在，那么对于当前主体来说仍然是：

\[
Tool
=
ExternalFunctionalResource.
\]

这里需要注意“本体主体”和“接口角色”的区别。一个 Tool 背后当然可能连接另一个真正独立的 Agent，但在当前 AgentOS 的进程模型中，我们处理的是接口关系：

\[
A
\leftrightarrow
A_{\mathrm{external}},
\]

而不是把 $A_{\mathrm{external}}$ 直接吸收成自身 Worker。此时它应通过明确的跨主体协议交互，其信息来源和权限关系也必须保持独立。

因此可以得到一个非常重要的边界：

\[
\boxed{
InternalWorker
\neq
ExternalAgent
}
\]

Internal Worker 的状态最终归属于主体：

\[
Owner(W_i)=A,
\]

而 External Agent 满足：

\[
Owner(A_j)=A_j.
\]

这一区分使 AgentOS 可以同时支持内部并发和真正的多 Agent 社会，而不混淆二者。内部 Worker 是“我的功能展开”；外部 Agent 是“另一个持续主体”。

从认知结构角度看，Worker 所承载的 Agent-CSO 最终仍然需要回流到主体生命周期。如果 Worker 在任务结束后发现重要结构 $C$，那么：

\[
AgentCSO_{W_i}(C)
\rightarrow
Integration_A(C)
\]

之后才能进入 $E$、$\Theta$ 或其他长期结构。否则 Worker 终止时，其局部发现就只是一次短暂计算，并未成为 Agent 的真正学习。

因此，“是否产生另一个 Agent”不能通过进程数量判断，而应该通过：

\begin{proposition*}
是否建立了独立生命周期因果链
\end{proposition*}

来判断。

\section{为什么主体原则上不可通过 fork、new 与 delete 任意复制}

传统计算机系统允许：

\[
fork(Process),
\]

甚至通过复制整个内存状态产生一个几乎完全相同的新进程。从信息状态角度看，我们似乎也可以对一个 Agent 做同样的事情：复制模型参数、工作记忆、情景记忆、结构记忆、自我状态和权限，然后得到两个相同实例。但是，只要 Agent 已经具有长期生命周期结构，这种操作会产生一个无法通过普通进程语义处理的问题：

\begin{statementbox}
复制之后，哪一个仍然是原来的持续主体？
\end{statementbox}

假设在 $t_0$ 时刻存在：

\[
A(t_0).
\]

如果执行：

\[
fork(A)
\rightarrow
A_1,A_2,
\]

并且：

\[
X_{A_1}(t_0)
=
X_{A_2}(t_0)
=
X_A(t_0),
\]

那么从 $t_0+\epsilon$ 开始，由于两个实例接受不同输入：

\[
X_{A_1}(t)
\neq
X_{A_2}(t).
\]

它们由此形成：

\[
\Gamma_{A_1}
\neq
\Gamma_{A_2}.
\]

也就是说，历史相同并不能保证未来主体唯一。fork 会把一条生命周期轨迹转换成两条后继轨迹：

\[
\Gamma_A
\rightarrow
\{
\Gamma_{A_1},
\Gamma_{A_2}
\}.
\]

如果系统仍然把二者都标记成同一 Agent，那么所有与：

权限；

责任；

长期目标；

记忆所有权；

身体控制权；

社会关系

相关的语义都会发生歧义。

因此在第一代 AgentOS 中，我们提出：

\[
\boxed{
fork(A)=Forbidden
}
\]

并进一步：

\[
\boxed{
new(A)=Forbidden,
\qquad
delete(A)=Forbidden
}
\]

这里需要非常谨慎地理解这些表达。它们不是说系统永远不能创建新的 Agent，而是说\textbf{一个已经存在的持续主体不能把创建、复制和删除主体自身当作普通运行时操作}。新 Agent 的创建应属于更高层的系统治理事件，而不是像创建线程一样由普通任务执行。

同理，“删除主体”也不能等价于：

\[
kill(pid).
\]

因为一个持续主体可能拥有长期责任、外部关系、资产、权限和历史状态。因此主体停止运行至少应区分：

\[
Sleep,
\quad
Shutdown,
\quad
Archive,
\quad
IrrecoverableTermination.
\]

其中前三者都可能保持：

\[
IdentityContinuity=1,
\]

而最后一种才真正终止生命周期轨迹。

因此：

\[
\boxed{
ExecutionTermination
\neq
IdentityTermination
}
\]

Checkpoint 恢复也同样不能被解释成创建新 Agent。如果主体在 $t_1$ 时刻崩溃，然后从 $t_0<t_1$ 的 checkpoint 恢复，那么我们面对的是：

\[
A(t_0)
\rightarrow
A'(t_1),
\]

而不是自动创建一个完全新的主体。系统需要显式记录“经历区间丢失”和“恢复事件”，将其纳入后续生命周期，而不能伪装成：

\[
t_0\rightarrow t_1
\]

之间从未发生异常。

这体现出一个更深的原则：

\begin{proposition*}
身份连续性依赖可解释的因果连续，
而不是状态向量的逐比特一致。
\end{proposition*}

因此，本章禁止普通主体 fork，真正保护的并不是某一块内存，而是生命周期因果关系的唯一性。

\section{身份连续性的系统级条件}

如果我们不允许简单通过语言模型的自我声明来定义“同一个 Agent”，那么身份连续性必须转化为可以被系统设计和验证的条件。一个持续 Agent 的身份并不要求所有内部变量恒定；恰恰相反，真正的持续 Agent 必须能够学习，因此：

\[
E(t+\Delta t)
\neq
E(t),
\]

\[
\Theta(t+\Delta t)
\neq
\Theta(t),
\]

甚至：

\[
\Psi(t+\Delta t)
\neq
\Psi(t).
\]

如果身份要求状态不变，那么学习本身就会破坏身份。因此身份连续必须定义成一种\textbf{受约束的历史连续性}。

可以将身份连续性定义为多个条件的联合：

\[
\mathcal{I}_A
=
F(
I_H,
I_M,
I_S,
I_C,
I_A,
I_B
),
\]

其中：

\[
I_H
=
\text{History Continuity},
\]

\[
I_M
=
\text{Memory Continuity},
\]

\[
I_S
=
\text{Self-Model Continuity},
\]

\[
I_C
=
\text{Control Continuity},
\]

\[
I_A
=
\text{Authority Continuity},
\]

\[
I_B
=
\text{Boundary Continuity}.
\]

历史连续性要求每次状态变化都有可追踪前驱：

\[
X(t+\Delta t)
=
F(
X(t),
O(t),
A(t),
L(t)
),
\]

而不能出现无来源的大规模状态替换。

记忆连续性并不要求所有经历永久保存，而要求新的长期结构能够在原则上追溯其形成历史。即使某些具体情景被压缩或遗忘，其主要结构变化仍具有可解释来源：

\[
E
\rightarrow
\Theta
\rightarrow
\Psi.
\]

自我连续性要求：

\[
\|\Psi(t+\Delta t)-\Psi(t)\|
\]

通常受到比普通工作记忆更严格的变化约束，即：

\[
\tau_\Psi
\gg
\tau_h.
\]

控制连续性则意味着主体级决策权不能未经治理突然迁移到局部 Worker。可以写成：

\[
\operatorname{RootAuthority}(t)
=
A
\]

在正常生命周期内保持不变。

权限连续性要求外部系统能够明确知道：

“当前行动由哪个持续主体授权？”

而边界连续性则要求 Agent 的身体、工具、数据和外部资源的归属变化经过显式授权，而不是随着任务上下文隐式漂移。

因此可以提出一个系统级身份连续条件：

\[
\boxed{
\mathcal I_A(t,t+\Delta t)
\geq
\theta_I
}
\]

其中 $\mathcal I_A$ 不是简单相似度，而是一种综合因果连续指标。

这里尤其要强调：

\[
\boxed{
IdentityContinuity
\neq
StateSimilarity
}
\]

两个状态非常相似的 Agent 可能是不同主体，例如 fork 后的两个实例；而同一个主体在经历多年学习以后，状态可能发生巨大变化，但仍然保持同一生命周期身份。

这与人类身份直觉也有一定相似性，但本章不依赖人类哲学人格理论。对于 AgentOS 来说，我们需要的是能够工程实现的最低条件：存在唯一的生命周期 ID、连续的事件日志、受治理的慢变量变化、唯一主体控制权、可验证 checkpoint 链以及明确的记忆归属。

因此，身份连续性最终不是一个单独模块，而是一组跨越：

\[
Runtime,
\quad
Memory,
\quad
Self,
\quad
Authority,
\quad
Recovery
\]

的系统不变量。

\section{单一主体原则的边界、价值与理论位置}

到这里，我们可以更加准确地理解为什么 AgentOS 要首先采用单一持续主体原则。它并不是因为“一切真正智能都必须只有一个中心”，也不是说未来绝对不可能存在分布式主体。它真正承担的是一个更有限、也更工程化的任务：在我们尚未解决“主体复制、分裂、合并和多主体身份归属”的完整理论之前，为持续人工智能建立一个可治理的最小主体模型。

这一模型可以压缩为：

\[
\boxed{
N_{\mathrm{Agent}}^{D_1}=1
}
\]

\[
\boxed{
|\mathcal T|\geq 0,
\qquad
|\mathcal W|\geq 0
}
\]

即一级主体唯一，而任务和 Worker 数量原则上不受该约束。

再进一步：

\[
\boxed{
Owner(T_i)=A
}
\]

\[
\boxed{
Owner(W_j)=A
}
\]

以及：

\[
\boxed{
D_2\nrightarrow Trap(D_1).
}
\]

这四组关系共同组成 AgentOS 的基础主体不变量。

从运行效率角度看，单一主体原则允许系统获得高度并发；从身份角度看，它防止并发退化成无意义的主体复制；从记忆角度看，它确保所有真正长期学习最终汇入同一条生命周期；从控制角度看，它保持 Root Authority 的唯一性；从安全角度看，它防止局部任务突然获得主体级权限；从软件架构角度看，它又使 AgentOS 能够使用成熟的任务调度、Worker Pool、异步消息和分布式执行技术，而不必将每个执行单元都包装成“人格”。

因此：

\[
\boxed{
SingleSubject
\neq
CentralizedComputation
}
\]

反而更准确的是：

\[
\boxed{
SingleSubject
+
DistributedFunctionalProcesses
}
\]

这也是本章最重要的系统思想之一。

从更深层的智能动力学角度看，主体其实是许多不同时间尺度闭环之间的一种长期因果闭合。Body 的快速闭环、工作记忆的秒级闭环、经验记忆的中慢尺度闭环、自我和生命周期生成结构的极慢闭环，并不是因为共享同一个 CPU 才成为一个主体，而是因为这些过程持续共同作用于同一条生命周期轨迹：

\[
\Gamma_A.
\]

因此可以给出本章最终定义：

\begin{definition*}
持续主体不是一个计算位置，而是一条在多尺度认知过程之间保持唯一因果归属、记忆归属和控制归属的生命周期动力学轨迹。
\end{definition*}

由此，AgentOS 的进程模型应当明确区分：

\[
\boxed{
Subject
}
\]

与

\[
\boxed{
Process.
}
\]

Process 是可创建、可替换、可暂停、可并行化的功能结构；Subject 则是这些结构长期服务和共同塑造的生命周期连续体。

这一区分也解释了为什么我们不能简单把传统操作系统中的 process 概念直接映射成 Agent。传统 OS 解决的是：

\[
\boxed{
WhoOwnsComputationalResources?
}
\]

而 AgentOS 还必须进一步解决：

\[
\boxed{
WhoOwnsTheExperience?
}
\]

\[
\boxed{
WhoOwnsTheMemory?
}
\]

\[
\boxed{
WhoOwnsTheDecision?
}
\]

以及最终：

\[
\boxed{
WhoContinuesToBecome?
}
\]

当这四个问题都指向同一个持续生命周期结构时，我们才真正拥有一个能够在任务变化、身体变化、记忆增长和环境变化中持续存在的 Agent。

因此，本章的核心并不是限制 AgentOS 的并发能力，而是在高度并发计算之上建立一个不会被任务、Worker、恢复和工具调用轻易撕裂的主体边界。只有在这一基础上，长期记忆、自我发展、权限成长以及更复杂的生命周期自治才具有稳定的系统承载对象。

本章由此确立：

\begin{statementbox}
Many Tasks,\ Many Processes,\ Many Loops,\ One Persistent Subject.
\end{statementbox}

这就是 AgentOS 第一代持续主体进程模型的基本形式。
